\section{新入寮生の皆さんへ　ー熊野寮で「予後」よく過ごすにはー}
\rightline{（文責：ラーメンマン）}

　新入寮生の皆さんこんにちは。

　皆さんが熊野寮に入寮する際、部屋を自分で選ぶことはできません。空いている部屋に、生活リズムが合うか等を可能な限り考慮して入れることにしていますが、近年の入寮者数の増加により、なかなかそれもうまくいっていないのが現状です。

　このような状況において、新入寮生の皆さんができるだけ「予後」よく寮生活を過ごすことができるよう、本記事では\textbf{「同部屋や、入寮した際に加わることになったコミュニティの人たちがしていたら注意するべき言動」}をご紹介します。熊野寮にはたくさんのコミュティがあるので、入るコミュニティや出会う人々によって寮生活の「質」が大きく変化する可能性があります。なお、本記事は特定の個人ないしはコミュニティを指して批判するものではありません。悪しからず。

\subsection{事例1：寮の基本的な会議に出ていないにもかかわらずしたり顔で寮の悪口を言う人}
　熊野寮は自治寮です。自治寮を自治寮たらしめているのは、会議における、議論による意思決定です。熊野寮では意思決定の根幹を定期的に実施される会議にて行っていますが、そのような会議に出ない（サボリ）にもかかわらず、したり顔で寮生や組織の方針への悪口（批判というレベルにも至らない）を吹聴する人がいます。こういった人に出会ったらば注意する必要があるでしょう。第一に、こういった人は会議に出ていない＝寮内の情報をきちんと把握できていない、ということなので、寮への解像度が低いです。第二に、寮生の議論の根幹である会議をサボって悪口を言うということは寮自治に対する責任を一切取る気がないということです。解像度の低い、責任を取る気もない言説をなるほどと受容してしまうのは、「だまされる」ことにも繋がります。もちろん会議に出ない人の中には研究の都合等で「出られない」人もいると思うので、そのあたりはよく見極めましょう。夜遅くまで起きているにもかかわらず会議にでなかったり、「寮における民主主義が～」などよくわからない言い訳をしているならば、まず意識的にサボっているとみて問題ないでしょう。

\subsection{事例2：大きな新歓や会議、寮食、イベントがあるにも関わらず部屋で鍋をしようとする人}
　熊野寮の4月は新歓シーズンです。毎日のように新歓があります。特に土日は大規模な新歓が行われます。また熊野寮には食堂があり、授業のある平日には寮食が提供されます。このような状態にも関わらず、自分の部屋やクローズドなコミュニティで鍋や炊き出しを行う人々がいます（ごくたまにではなく、頻繁に）。新歓は寮生同士が知り合い、人脈を広げていく大変重要な機会です。また、寮食を喫食することは寮自治と福利厚生の面から非常に重要な熊野寮食堂を防衛していくために寮生がするべきことです（入寮する際にきちんと説明される部分です）。新入寮生を重要な交流の機会や寮生の務めから引きはがし、特定少数のコミュニティにとどめおこうとする人たちは、寮のことや新入寮生のことを考えていません。寮生が様々な形で力を合わせて寮自治を行っていることを忘れ、自分たちの権益を確保し、再生産することを考えているにすぎません。皆さんが新歓やイベントに出たり、寮食を喫食する中でたくさんの寮生と交流することを恐れているのです。

　せっかく熊野寮に入ったわけですので、寮生としての務めを果たしながら、多くの人々と交流し、実りある寮生活を送る事が望ましいだろうと考えます。自分本位・自分たち本位の集団は周りから困った人たちだととらえられ、深さや豊かさのある寮生活からは程遠くなるでしょう。寮祭などの楽しいイベントも100\%楽しむことができなくなるかもしれません。具体的な事例を紹介しましたので、これらを念頭に少しでも「予後」の良い寮生活を送ることができるようお祈りします。